\section{Introduction}

Large Language Models (LLMs) are increasingly used to process unstructured data due to their strong capabilities in semantic understanding and reasoning.
Modern data platforms have begun to develop LLM-powered semantic operators or User-Defined Functions (UDFs), including commercial systems such as Google BigQuery~\cite{google2026bigqueryai}, Snowflake Cortex AI~\cite{snowflake2026cortexai}, and Alibaba PolarDB for AI~\cite{alibaba2026polardbai}, as well as prototype systems for semantic query processing~\cite{DBLP:journals/pvldb/PatelJPGAGZ25}.
This integration poses new problems for physical optimization in database systems. 
Unlike classical relational operators, LLM-powered operators invoke a model on individual records or batches. Each request repeats the ranking instruction and includes the text of every row it compares. Reducing the request count can therefore save tokens, but a plan with fewer requests can still cost more if it resubmits more text.

Among semantic operators, semantic \textsf{ORDER BY} supports queries whose ordering criterion is expressed in natural language.
For example, ordering reviews by their helpfulness for a purchase decision, passages by the strength of their evidential support, or retrieved documents by their relevance to a user question. 
Such functionality is central to document reranking in retrieval-augmented generation (RAG)~\cite{DBLP:conf/nips/YuP0WYZSC24} and semantic search. As semantic \textsf{ORDER BY} requires an LLM to infer a global ordering over input items, when the input exceeds the effective context window of the LLM, no single invocation can inspect all items and produce a global ordering. 
A natural solution is divide-and-conquer sorting: inferring the local order of item subsets that fit in a context window, and combining local orders into a global order. 
This makes processing Semantic \textsf{ORDER BY} substantially more difficult than processing record-independent operators such as semantic filtering. 


Existing LLM ranking methods use pointwise scores, pairwise comparisons, or listwise permutations. Pointwise methods score items independently, which makes scores from different prompts difficult to calibrate~\cite{DBLP:conf/emnlp/DrozdovZD0RWAIM23,DBLP:conf/emnlp/SachanLJAYPZ22}. Pairwise methods compare two items at a time and combine these judgments into a ranking~\cite{DBLP:conf/acl/LuoCH024,DBLP:conf/naacl/QinJHZWYSLLMWB24}. We focus on listwise ranking, which orders an entire batch in one call~\cite{DBLP:conf/emnlp/0001YMWRCYR23,DBLP:journals/corr/abs-2603-17223,DBLP:journals/corr/abs-2509-00303}. The physical operator decides how to turn these local permutations into a full order or an ordered prefix.

Existing operators organize these calls in different ways. ListK studies tournament selection and multi-pivot partitioning~\cite{DBLP:journals/corr/abs-2603-17223}. TBS retains prior judgments in a preference graph and skips batches whose pairwise orders are already implied~\cite{akinola2026tbs}. Two-way external merge instead generates sorted runs and uses their ordered heads to assemble the result~\cite{DBLP:journals/corr/abs-2509-00303}. A merge call can contain rows whose relative order is already known. A wider fan-in reduces the number of passes, but it also changes how many rows can enter each call and how often they return to the prompt. The number of known pairs in a batch does not tell us which plan spends fewer input tokens.

The key observation is that a run needs to supply only enough rows to determine what can be emitted next. A sliding window retains rows that cannot yet be emitted and replenishes each active run. A segmented loser tree retains the losing segments and reranks only one path on a refill. These two choices of merge group expose different trade-offs between the number of calls and the number of row appearances. Output demand changes the trade-off again: Full ranking drains every group, while a small \LIMIT~$K$ can stop after consuming uneven prefixes from different children, even though both modes generate all initial runs.

Figure~\ref{fig:framework} illustrates why one plan need not serve both queries. Given the same initial runs, the planner can choose different mergers, fan-ins, and segment lengths for Full and \LIMIT~$K$. It estimates the work of each feasible plan under the requested output demand and prices it using the query's prompt profile.

\textbf{Contributions.}
\begin{itemize}
  \item \textbf{Multi-way listwise merge.} We develop a sliding-window merger with safe-prefix emission and a segmented loser tree with pointer-guided refresh. The same groups support materialized full ranking and demand-driven ordered prefixes. We establish correctness under a consistent ordering oracle.
  \item \textbf{Demand-aware plan selection.} We estimate the calls and row appearances of each feasible merge plan by replaying its finite execution. A token profile fits the fixed request cost and per-row cost without labeled rankings or model generation. The planner chooses the family, fan-in, and quota or segment length for the requested demand.
  \item \textbf{Evaluation.} We study end-to-end ranking quality, token use, and latency, and isolate plan selection, input scale, prefix caching, and the change from small limits to full ranking.
\end{itemize}
